% ===================================================================
% AB-05d: Mi tiempo libre - Wiederholung
% Spanisch Klasse 7 - Sequenz 5: Mi tiempo libre
% ===================================================================
% Abgeleitet von: LE-05d.tex (Score: 9.2/10)
% Punkte: 20 (4+4+6+6)
% ===================================================================

\documentclass[11pt, a4paper]{article}

% ===================================================================
% SW-MODUS TOGGLE
% ===================================================================
\newif\ifswmodus
\swmodusfalse

% ===================================================================
% PAKETE
% ===================================================================

\usepackage[utf8]{inputenc}
\usepackage[T1]{fontenc}
\usepackage[ngerman]{babel}
\usepackage{helvet}
\renewcommand{\familydefault}{\sfdefault}

\usepackage[margin=2cm]{geometry}
\usepackage{enumitem}
\usepackage{tabularx}
\usepackage{booktabs}
\usepackage{multirow}

\usepackage{xcolor}
\usepackage{framed}
\usepackage{tcolorbox}
\tcbuselibrary{skins}

\usepackage{fancyhdr}
\usepackage{lastpage}
\usepackage{needspace}

% ===================================================================
% FARBEN
% ===================================================================

\definecolor{spanisch-color}{HTML}{c41e3a}
\definecolor{grau-color}{HTML}{888888}
\definecolor{hellgrau-color}{HTML}{f5f5f5}
\definecolor{orange-color}{HTML}{b35c00}

\definecolor{sw-dark}{HTML}{000000}
\definecolor{sw-gray}{HTML}{666666}
\definecolor{sw-white}{HTML}{ffffff}

\ifswmodus
    \colorlet{spanisch}{sw-dark}
    \colorlet{grau}{sw-gray}
    \colorlet{hellgrau}{sw-white}
    \colorlet{orange}{sw-dark}
\else
    \colorlet{spanisch}{spanisch-color}
    \colorlet{grau}{grau-color}
    \colorlet{hellgrau}{hellgrau-color}
    \colorlet{orange}{orange-color}
\fi

\newcommand{\fachfarbe}{spanisch}

% ===================================================================
% VARIABLEN
% ===================================================================

\newcommand{\thema}{Mi tiempo libre}
\newcommand{\sequenz}{05}
\newcommand{\block}{d}
\newcommand{\fach}{Spanisch}
\newcommand{\klasse}{7}
\newcommand{\maxpunkte}{20}
\newcommand{\datum}{\today}

% ===================================================================
% KOPF- UND FUSSZEILE
% ===================================================================

\pagestyle{fancy}
\fancyhf{}
\renewcommand{\headrulewidth}{0.4pt}
\renewcommand{\footrulewidth}{0.4pt}

\lhead{\fach{} -- Klasse \klasse}
\chead{AB-\sequenz\block}
\rhead{\datum}

\lfoot{}
\cfoot{}
\rfoot{Seite \thepage\ von \pageref{LastPage}}

% ===================================================================
% BEFEHLE
% ===================================================================

\newcommand{\punktebox}[1]{\hfill\fbox{\small \_/#1}}
\newcommand{\luecke}[1]{\rule{#1}{0.4pt}}

\newtcolorbox{merkkasten}[1][]{
    colback=hellgrau,
    colframe=\fachfarbe,
    colbacktitle=white,
    coltitle=\fachfarbe,
    boxrule=1pt,
    arc=0pt,
    left=8pt, right=8pt, top=8pt, bottom=8pt,
    fonttitle=\bfseries,
    attach boxed title to top left={yshift=-2mm, xshift=5mm},
    boxed title style={boxrule=0pt, colback=white},
    title=#1
}

\newtcolorbox{vokabelkasten}{
    colback=white,
    colframe=\fachfarbe,
    boxrule=0.5pt,
    arc=0pt,
    left=5pt, right=5pt, top=5pt, bottom=5pt
}

\newtcolorbox{hinweis}{
    colback=white,
    colframe=orange,
    colbacktitle=white,
    coltitle=orange,
    boxrule=1pt,
    arc=0pt,
    left=5pt, right=5pt, top=5pt, bottom=5pt,
    fonttitle=\bfseries,
    attach boxed title to top left={yshift=-2mm, xshift=5mm},
    boxed title style={boxrule=0pt, colback=white},
    title=Hinweis
}

\newenvironment{aufgabe}[1]{%
    \needspace{5\baselineskip}%
    \nopagebreak[4]%
    \vspace{10pt}
    \noindent\textbf{\textcolor{\fachfarbe}{Aufgabe #1}}
    \nopagebreak[4]%
    \vspace{5pt}
    \nopagebreak[4]%
    \begin{list}{}{\leftmargin=0pt}
    \item[]
}{%
    \end{list}
}

\newcommand{\es}[1]{\textcolor{\fachfarbe}{\textbf{#1}}}

% ===================================================================
% DOKUMENT
% ===================================================================

\begin{document}

% ===================================================================
% KOPFBEREICH
% ===================================================================

\begin{center}
    {\Large\bfseries\textcolor{\fachfarbe}{Arbeitsblatt: \thema{} -- Wiederholung}}
\end{center}

\vspace{5pt}

\noindent
\begin{tabularx}{\textwidth}{|X|X|X|}
\hline
\textbf{Name:} \luecke{4cm} &
\textbf{Klasse:} \klasse &
\textbf{Datum:} \luecke{2cm} \\
\hline
\end{tabularx}

\vspace{15pt}

% ===================================================================
% MERKKASTEN: Verben auf -er
% ===================================================================

\begin{merkkasten}[Merkkasten 1: Verben auf -er]

\begin{tabular}{l|cccccc}
& \textbf{yo} & \textbf{tú} & \textbf{él/ella} & \textbf{nosotros} & \textbf{vosotros} & \textbf{ellos} \\
\hline
\es{comer} & \luecke{1.5cm} & \luecke{1.5cm} & \luecke{1.5cm} & \luecke{1.5cm} & coméis & \luecke{1.5cm} \\
\end{tabular}

\vspace{0.5em}
\textbf{Weitere -er Verben:} \es{beber} (trinken), \es{leer} (lesen), \es{correr} (laufen)

\end{merkkasten}

\vspace{10pt}

% ===================================================================
% MERKKASTEN: Verben auf -ir
% ===================================================================

\begin{merkkasten}[Merkkasten 2: Verben auf -ir]

\begin{tabular}{l|cccccc}
& \textbf{yo} & \textbf{tú} & \textbf{él/ella} & \textbf{nosotros} & \textbf{vosotros} & \textbf{ellos} \\
\hline
\es{vivir} & \luecke{1.5cm} & \luecke{1.5cm} & \luecke{1.5cm} & \luecke{1.5cm} & vivís & \luecke{1.5cm} \\
\end{tabular}

\vspace{0.5em}
\textbf{Weitere -ir Verben:} \es{escribir} (schreiben), \es{abrir} (öffnen)

\end{merkkasten}

\vspace{10pt}

% ===================================================================
% AUFGABE 1: -er Verben konjugieren
% ===================================================================

\begin{aufgabe}{1 -- Konjugiere die Verben auf -er}
Ergänze die richtige Form. \punktebox{4}
\end{aufgabe}

\noindent
a) Yo \luecke{2.5cm} (comer) una pizza. \hfill
b) Tú \luecke{2.5cm} (beber) agua. \\[0.8em]
c) Ella \luecke{2.5cm} (leer) un libro. \hfill
d) Nosotros \luecke{2.5cm} (comer) en casa. \\

\vspace{10pt}

% ===================================================================
% AUFGABE 2: -ir Verben konjugieren
% ===================================================================

\begin{aufgabe}{2 -- Konjugiere die Verben auf -ir}
Ergänze die richtige Form. \punktebox{4}
\end{aufgabe}

\noindent
a) Yo \luecke{2.5cm} (vivir) en Rostock. \hfill
b) Tú \luecke{2.5cm} (escribir) una carta. \\[0.8em]
c) Él \luecke{2.5cm} (vivir) en España. \hfill
d) Nosotros \luecke{2.5cm} (escribir) en el cuaderno. \\

\vspace{10pt}

% ===================================================================
% MERKKASTEN: gustar
% ===================================================================

\begin{merkkasten}[Merkkasten 3: gustar + Infinitiv]

\begin{tabular}{ll}
\es{Me gusta} + Infinitiv & = Mir gefällt es zu... / Ich mag... \\
\es{Te gusta} + Infinitiv & = Dir gefällt es zu... / Du magst... \\
\es{Le gusta} + Infinitiv & = Ihm/Ihr gefällt es zu... \\
\end{tabular}

\vspace{0.5em}
\textbf{Beispiel:} \es{Me gusta jugar al fútbol.} = Ich spiele gern Fußball.

\end{merkkasten}

\vspace{10pt}

% ===================================================================
% AUFGABE 3: gustar-Sätze
% ===================================================================

\begin{aufgabe}{3 -- Bilde Sätze mit gustar}
Schreibe vollständige Sätze. \punktebox{6}
\end{aufgabe}

\noindent
a) yo / jugar al fútbol

\vspace{0.5em}
\noindent
\luecke{12cm}

\vspace{0.8em}

\noindent
b) tú / escuchar música

\vspace{0.5em}
\noindent
\luecke{12cm}

\vspace{0.8em}

\noindent
c) él / leer libros

\vspace{0.5em}
\noindent
\luecke{12cm}

\vspace{10pt}

% ===================================================================
% AUFGABE 4: Hobbys zuordnen
% ===================================================================

\begin{aufgabe}{4 -- Ordne die Hobbys zu}
Welches Verb passt? Verbinde mit \es{jugar}, \es{tocar} oder \es{hacer}. \punktebox{6}
\end{aufgabe}

\begin{center}
\begin{tabular}{rcl}
al fútbol & \luecke{2cm} & \multirow{3}{*}{\es{jugar}} \\
al baloncesto & \luecke{2cm} & \\
a videojuegos & \luecke{2cm} & \\[0.8em]
la guitarra & \luecke{2cm} & \multirow{2}{*}{\es{tocar}} \\
el piano & \luecke{2cm} & \\[0.8em]
deporte & \luecke{2cm} & \multirow{2}{*}{\es{hacer}} \\
los deberes & \luecke{2cm} & \\
\end{tabular}
\end{center}

% ===================================================================
% PUNKTEÜBERSICHT
% ===================================================================

\vspace{20pt}
\hrule
\vspace{10pt}

\begin{flushright}
\textbf{Punkte:} \luecke{1cm} / \maxpunkte
\end{flushright}

\end{document}
